\section{Compact embedding and the five-dimensional gauging matrix}
\label{sec:compact-gauging}

The physical matrix correlating the coordinates of
Section~\ref{sec:local-sectors} has two types of columns. Exceptional curves at
nodes give electrically charged wrapped-M2 hypermultiplets. Roots in an
exceptional del Pezzo surface determine how compact Abelian vectors embed into
the flavor Cartan of the interacting fixed point. Both columns are obtained by
evaluating curve classes on compact divisors.

\subsection{Vector fields and wrapped-M2 charges}
\label{subsec:m2-charges}

Let $X$ be a compact singular Calabi--Yau threefold and let
\begin{equation}
 \pi:\wideX\longrightarrow X
 \label{eq:crepant-resolution}
\end{equation}
be a projective crepant resolution. Choose harmonic two-forms $\omega_I$ on
$\wideX$ whose cohomology classes span the vector directions relevant to the
transition. The M-theory three-form has the expansion
\begin{equation}
 C_3=\sum_I A^I\wedge\omega_I+\cdots .
 \label{eq:c3-reduction}
\end{equation}
If an M2-brane wraps a curve $C\subset\wideX$, its worldline coupling is
\begin{equation}
 \int_{\mathbb R\times C}C_3
 =\sum_I\left(\int_C\omega_I\right)\int_{\mathbb R}A^I.
 \label{eq:m2-coupling}
\end{equation}
For a divisor class $D_I$ dual to $\omega_I$, the electric charge is therefore
\begin{equation}
 c_I(C)=\int_C\omega_I=D_I\cdot C.
 \label{eq:m2-charge}
\end{equation}

At a resolved node $a$, let $C_a\cong\mathbb P^1$ denote the exceptional
curve. The wrapped M2-brane is the massless hypermultiplet of the conifold
transition, with charge vector
\begin{equation}
 c_{Ia}=D_I\cdot C_a.
 \label{eq:nodal-charge}
\end{equation}
Thus the map
\begin{equation}
 \ZZ^{N_{\mathrm{node}}}\longrightarrow
 H_2(\wideX;\ZZ)/\mathrm{tors},
 \qquad e_a\longmapsto[C_a]
 \label{eq:nodal-homology-map}
\end{equation}
is simultaneously the homology map of the resolution and the electric charge
map of the transition states. This is the standard compact conifold
description \cite{Katz:1996,Mohaupt:2004de}.

\subsection{Flavor cocharacters from del Pezzo roots}
\label{subsec:flavor-cocharacters}

Suppose the exceptional locus over $p$ contains
$S_{8,p}=\PP^1\times\PP^1$, $S_{7,p}$, $S_{6,p}$ or $S_{5,p}$.
The local non-Abelian flavor lattice is the root
lattice inside $K_{S_p}^{\perp}$, the orthogonal complement of $K_{S_p}$ in
$H_2(S_p;\ZZ)$ for the intersection form. Let
\begin{equation}
 \iota_p:S_p\hookrightarrow\wideX
\end{equation}
be the inclusion and let $R_{p\alpha}$ be a root. The restriction
$D_I|_{S_p}$ assigns degrees to the M2 states supported on curves in $S_p$.
For $E_2$, with $R=e_1-e_2$, the difference of the charges of the two
exceptional-curve states is
\begin{equation}
 c_I(e_1)-c_I(e_2)
 =D_I\cdot\iota_{p*}(e_1-e_2).
 \label{eq:e2-charge-difference}
\end{equation}
This is the differential of the homomorphism from the compact $U(1)$ generated
by $A^I$ into the $SU(2)$ flavor torus. More generally, define
\begin{equation}
 m_{Ip\alpha}
 =\left\langle D_I|_{S_p},R_{p\alpha}^{\vee}\right\rangle
 \sim D_I\cdot\iota_{p*}R_{p\alpha}.
 \label{eq:cocharacter}
\end{equation}
The proportionality allows for the conventional common normalization between
a primitive geometric root and a Cartan generator. All algebras here are
simply laced. We fix compatible root-marked moment-map coordinates throughout.
Changing a root/coroot normalization or orientation changes the corresponding
coordinate and matrix column together. The resulting kernels are identified
by this invertible coordinate change; their coordinate vectors are not
literally unchanged.

For $E_2$ there is one reduced root coordinate $\beta_p$. For $E_3$, the
$A_2$ component uses $(R_{p1},R_{p2})$ and the $A_1$ component uses $R_{p0}$,
as in \eqref{eq:e3-roots}. The $E_1$ coordinate $u_s$ multiplies the ruling
difference, and the $E_4$ coordinates $c_{vj}$ multiply the four roots in
\eqref{eq:e4-simple-roots}. Gauging a flavor isometry couples the
vector-multiplet auxiliary triplet to the corresponding current-multiplet
moment map \cite{Ceresole:2000}. Since moment maps add under diagonal
gauging, the compact complex moment maps are
\begin{equation}
 \begin{aligned}
 \mu_I^{\CC}
 &=\sum_a c_{Ia}t_a+\sum_s m^{(8)}_{Is}u_s
  +\sum_r m^{(7)}_{Ir}\beta_r\\
 &\quad+\sum_p\left(
     m_{Ip1}H_{p1}+m_{Ip2}H_{p2}+m_{Ip0}L_p
   \right)+\sum_v\sum_{j=1}^4m^{(5)}_{Ivj}c_{vj}.
 \end{aligned}
 \label{eq:compact-complex-mm}
\end{equation}
This conclusion uses only the flavor current multiplets of the interacting
SCFT; it does not require a weakly coupled Lagrangian description of the
$E_1$, $E_2$, $E_3$ or $E_4$ fixed point.

\subsection{Branch profiles and the compact charge map}
\label{subsec:branch-profiles}

At each $E_3$ point choose either the full Higgs component
$\cH_{\mathfrak{sl}_2}$ or $\cH_{\mathfrak{sl}_3}$. A collection of choices
is a \emph{branch profile} $\sigma$. The associated local coordinate space
uses the invariant images of \eqref{eq:e3-invariant-images}:
\begin{equation}
 \begin{aligned}
 \cB_{\sigma}
 &=\bigoplus_a\CC_{t_a}\oplus\bigoplus_s\CC_{u_s}
  \oplus\bigoplus_r\CC_{\beta_r}
  \oplus\bigoplus_p \cB_{p,\sigma_p}
  \oplus\bigoplus_v\CC^4_{c_{v1},\ldots,c_{v4}},\\
 \cB_{p,\mathfrak{sl}_2}&=\CC_{L_p},
 \quad
 \cB_{p,\mathfrak{sl}_3}=\CC^2_{H_{p1},H_{p2}}.
 \end{aligned}
 \label{eq:profile-space}
\end{equation}
Let $\Lambda_{\mathrm{vec}}$ be the lattice generated by the compact
Abelian vector directions $A^I$, equivalently by their divisor classes
$D_I$ modulo the integral principal-divisor relations; the $D_I$ are taken to
span $H^2(\wideX;\QQ)$. Under the hypotheses of
Theorem~\ref{thm:first-order-image}, the singularities are rational and
Serre duality gives $H^2(\wideX,\cO_{\wideX})=0$. Projectivity and the
Lefschetz $(1,1)$ theorem therefore ensure that divisor classes span
$H^2(\wideX;\QQ)$. No ambient toric presentation is required. On the admissible
class of Section~\ref{subsec:admissible-class}, restricted ambient toric
divisors suffice. The dual lattice records one complex moment-map equation for each vector
direction. The restriction of
\eqref{eq:compact-complex-mm} to a profile is therefore the linear map
\begin{equation}
 \Qcompact:\cB_{\sigma}\longrightarrow
 \Lambda_{\mathrm{vec}}^{\vee}\otimes\CC.
 \label{eq:compact-charge-map}
\end{equation}

The same map can be defined before complexification. Let
$\Lambda^{\mathrm H}_{p,\sigma_p}$ be the root sublattice selected by the
local component, and let a nodal summand be generated by $[C_a]$. Then
\begin{equation}
 \begin{split}
 \Lambda^{\mathrm H}_{\sigma}
 &=\bigoplus_a\ZZ[C_a]
   \oplus\bigoplus_p\Lambda^{\mathrm H}_{p,\sigma_p},\\
 Q^{\mathrm{geo}}_{X,\sigma}:\Lambda^{\mathrm H}_{\sigma}
 &\longrightarrow H_2(\wideX;\ZZ)/\mathrm{tors},
 \qquad
 \gamma\longmapsto\sum_p\iota_{p*}\gamma_p .
 \end{split}
 \label{eq:geometric-charge-map}
\end{equation}
Evaluating \eqref{eq:geometric-charge-map} on the divisor classes gives the
matrix in \eqref{eq:compact-charge-map}. In other words, the geometric
local-to-global map and the physical embedding of the local electric and
flavor charges are dual descriptions of one lattice homomorphism.

In the toric hypersurface examples, the finite matrix is initially computed
using ambient divisor classes. The geometric comparison reviewed in
Section~\ref{sec:global-theorems} proves that, under the admissibility
hypotheses used here, its branch-restricted kernel equals the complete
topological kernel. Thus the finite presentation does not miss a continuous
constraint. Questions about a possible finite index between the printed toric
lattice and the full integral cohomology lattice affect global forms and line
operators, but not the rational kernel or the continuous Higgs selection rule.

\subsection{The projected moment-map locus}
\label{subsec:projected-locus}

Let $G_{\mathrm{eff}}$ be the compact Abelian group acting effectively on the
transition sectors through \eqref{eq:compact-complex-mm}. In the rigid
zero-level description, form the triplet-moment-map quotient of the product of
the local Higgs cones on a fixed profile. We use
$\cM_{H,\sigma}^{\mathrm{rigid}}$ for this quotient and retain only its map to
the invariant coordinates \eqref{eq:profile-space}.

\begin{proposition}[physical kernel]
\label{prop:physical-kernel}
For every branch profile $\sigma$, the image of the rigid zero-level Higgs
locus in its local deformation coordinates is
\begin{equation}
 \operatorname{Im}\!\left(
 \cM_{H,\sigma}^{\mathrm{rigid}}\longrightarrow\cB_{\sigma}
 \right)
 =\ker(\Qcompact).
 \label{eq:physical-kernel}
\end{equation}
\end{proposition}

\begin{proof}
If a point solves the complex moment-map equations, equation
\eqref{eq:compact-complex-mm} states exactly that its invariant coordinate
vector $b\in\cB_{\sigma}$ obeys $\Qcompact b=0$. This proves one inclusion.

Conversely, take $b\in\ker\Qcompact$. Lift every nodal coordinate $t_a$ by
\eqref{eq:conifold-balanced}. Lift every $E_2$ coordinate $\beta_r$ by
\eqref{eq:e2-double-cover} with equal absolute values. On an $E_3$ component,
and at every $E_1$ or $E_4$ point, use the balanced rank-one ADHM
representatives of Appendix~\ref{app:local-algebra}. Every local real Cartan moment map then vanishes
separately. The total real moment maps consequently vanish, while the total
complex moment maps vanish because $\Qcompact b=0$. This constructs a point
of the zero-level triplet-moment-map locus above every $b$ in the complex
kernel. Since the coordinates in $\cB_{\sigma}$ are gauge invariant, taking
the quotient does not change the image.
\end{proof}

If the selected local components have total quaternionic dimension
$d_{\mathrm{loc}}$ and a rank-$s$ Abelian group acts locally freely, the
familiar count gives
\begin{equation}
 \dim_{\HH}\cM_H=d_{\mathrm{loc}}-s
 =\dim_{\CC}\ker\Qcompact.
 \label{eq:dimension-check}
\end{equation}
Equation \eqref{eq:physical-kernel} is stronger than this count: it specifies
which combinations of coordinates survive and remains meaningful at profiles
where the action has nontrivial stabilizers. It also avoids the incorrect
identification of the full Higgs branch with a complex vector space.

\subsection{Compact supergravity qualification}
\label{subsec:supergravity-qualification}

The microscopic charge assignment \eqref{eq:m2-charge} and flavor embedding
\eqref{eq:cocharacter} are exact compactification data. Proposition
\ref{prop:physical-kernel} is stated in the rigid hyperk\"ahler-cone
description at zero moment-map level. The scalar manifold of five-dimensional
$\mathcal N=1$ supergravity is instead quaternionic K\"ahler, and its detailed
metric receives gravitational and strong-coupling corrections. We use the
rigid statement only to identify the holomorphic Higgs-deformation projection
and its compact gauging constraint. We do not claim that it constructs the
complete finite-Planck-mass scalar manifold.

This distinction does not turn the matrix into a background flavor constraint.
At finite compact volume the fields $A^I$ in \eqref{eq:c3-reduction} have
finite kinetic terms and are dynamical. Their auxiliary-field equations impose
the compact moment maps. What changes in a noncompact limit is the set of rows
that remain dynamical, as discussed in Section~\ref{sec:decoupling}.
